\chapter{Evaluation der Bechmark}\label{chap:conclusion}
In diesem Kapitel werden die Ergebnisse der begleitenden Benchmark (siehe \chapref{Links}) präsentiert und erörtert. 
% Brief wrapup of the achievements of this thesis.
% Directions for future work (if applicable).

% Experiments often include benchmarks.
% You need to describe you benchmarking setup: parameters of the machine used (architecture, memory, etc), parameters of the software used (OS version, libraries used and their version), the benchmark programs.
% The programs need not be part of the thesis, but they should be described and available separately.

% If you measure run times and evaluate your findings statistically, then there is great potential for errors.
% Consult SIGPLAN's checklist for empirical evaluations \url{http://www.sigplan.org/Resources/EmpiricalEvaluation/} for advice. 

\section{Aufbau der Benchmark}
Die hier benutzte Benchmark ist in python für das Framework pytest geschrieben, und untersucht sechs verschiedene Basis-Konfigurationen. Diese Konfigurationen unterscheiden sich in der Gesamtzahl der Testcases, die ausgeführt werden, sowie der Anzahl an Iterationen, die je Testcase ausgeführt werden, was die Ausführungsdauer eines Testcase bestimmt. Jede Konfiguration wird mit sechs verschieden gefüllten \texttt{archives} getestet.\\
Dabei werden die sogenannten parametrisierten Tests von pytest genutzt, dargestellt in \lstref{paratestPytest}:
\begin{lstlisting}[caption={Inhalt von test\_simple.py aus der Konfiguration ''bench\_fast\_few\_00''}, label={paratestPytest}]
@pytest.mark.parametrize("case", range(10))
def test_simple(case):
    index = 0
    for i in range(100):
        index += 1
    assert index == 100
\end{lstlisting}
range(x) in der obersten Zeile bestimmt, wie viele Kopien des Testcase mit durchnummerierten Namen beim Ausführen des Testcase erstellt werden.\\
range(y) in der for-Schleife in der vierten Zeile bestimmt die Anzahl der Iterationen, die im Testcase ausgeführt werden, bevor er terminiert.\\
Das in \lstref{paratestPytest} demonstrierte Beispiel erzeugt also 10 Testcases, welche jeweils 100 Iterationen ausführen.
Die Konfigurationen werden in 3 Geschwindigkeiten unterteilt: \textit{fast}, \textit{medium} und \textit{slow}. Für jede dieser Geschwindigkeiten existieren zwei Untergruppen: \textit{few} bezeichnet eine Konfiguration mit wenigen Testcases, \textit{many} eine mit vielen. Dabei wird die Gesamtmenge der Iterationen zur Erhaltung der Vergleichbarkeit innerhalb einer Geschwindigkeitsklasse gleich gehalten. Ein Beispiel: Während die Konfiguration ''fast\_few'' 10 Testcases mit je 100 Iterationen enthält, enthält die Konfiguration ''fast\_many'' 100 Testcases mit je 10 Iterationen. Somit werden in beiden Konfigurationen $10 \times 100=1000$ Iterationen ausgeführt, jedoch in unterschiedlichen Strukturen.\\
Jede Konfiguration enthält außerdem die Datei out.xml im Ordner archive. Diese Datei enthält das Ergebnis von vorherigen Testläufen, die vom Verfahren genutzt werden können, und erfüllt damit die Rolle der \texttt{archive}-Komponente (siehe \chapref{chap:Komponente_archive}). Dabei existieren für jede der sechs bisher vorgestelle Kombination aus Testcaseanzahl und Iterationsanzahl sechs Konfigurationen: Eine Konfiguration enthält ein leeres \texttt{archive}, folgend gibt es Konfigurationen mit \texttt{archives}, die jeweils 20, 40, 60, 80 und abschießend 100 Prozent der Testcases in der Konfiguration abdecken.

Durch die drei Geschwindigkeitsklassen, die beiden unterschiedlichen Mengen von Testcases, und den unterschiedliche gefüllten \texttt{archives} beträgt die Gesamtzahl der in der Benchmark überprüften Konfigurationen 36. 

\subsection{Spezifikationen des ausführenden Systems}
Der Computer, auf dem die Benchmark ausgeführt wurde, weist folgende Hardware-Spezifikationen auf:
\begin{itemize}
    \item Prozessor: AMD Ryzen 5 7600, 3801 MHz, 6 Kerne, 12 logische Prozessoren
    \item RAM: 32,0 GB DDR5-6000
    \item Speicher: 1TB NVMe SSD
\end{itemize}
Die Benchmark wurde unter wsl ausgeführt. Nähere Informationen zur Softwareumgebung sind aus der Datei flake.lock im Repository der Benchmark (siehe \chapref{Links}) zu entnehmen.

\section{Datenerhebung der Benchmark}
Es gilt zu beachten, dass die Benchmark die Dauer des gesamten Prozesses der Testerstellung misst. Für das in dieser Bacheloarbeit vorgestellte Verfahren bedeutet dies, dass die Ausführung der \texttt{testDiscovery}, der Ablauf von \texttt{testAuditor} sowie der Ablauf der \texttt{testExecution} gemessen werden. Die \texttt{testDiscovery} sowie die \texttt{testExecution} sorgen dabei für einen Testframeworkabhängigen Overhead. 

Um einen direkten Verleich der Ergebnisse zu ermöglichen, werden für jeden Durchlauf einer Konfiguration zwei idividuelle Durchläufe gestartet:\\
Der erste Lauf besteht auf der konventionellen Ausführung von Tests, bei der alle vorhandenen Testcases neu ausgeführt werden. Der zweite Lauf benutzt das in dieser Bachelorarbeit entwickelte Verfahren, um auf Grundlage der im \texttt{archive} liegenden Testausgaben die Mehrfache Ausführung von Testcases zu vermeiden.

Um die Auswirkungen von äußeren Einflüssen aus dem System auf den Ablauf der Benchmark bei der Datenerhebung auf die Ergebnisse zu verringern, wird die gesamte Benchmark zehnmal vollständig hintereinander ausgeführt. Das durchschnittliche Ergebnis ist dann die Grundage für \chapref{chap:result_benchmark}. 

\newpage
\section{Ergebnisse der Benchmark}\label{chap:result_benchmark}
$T$ bezeichnet die Anzahl an Testcases in der Konfiguration.\\
$I$ bezeichnet die Anzahl an Iterationen pro Testcase.
\begin{figure}[htbp]
    \centering
    \subfloat[wenige kurze Testcases: $T$=10, $I$=100]{%
        \resizebox{0.45\textwidth}{!}{%
        \benchmarkplot{
            (0, 181.7)
            (20, 159.4)
            (40, 160.1)
            (60, 160.1)
            (80, 160.1)
            (100, 159.6)}
            {
                (0, 305.6)
                (20, 300)
                (40, 295.7)
                (60, 291.9)
                (80, 290.0)
                (100, 161.8)
            }%
        }%
    }
    \hfill
    \subfloat[viele kurze Testcases: $T$=100, $I$=10]{%
        \resizebox{0.45\textwidth}{!}{%
        \benchmarkplot{
            (0, 196.7)
            (20, 197.2)
            (40, 193.0)
            (60, 195.0)
            (80, 194.0)
            (100, 194.2)
        }{
            (0, 377.7)
            (20, 363.9)
            (40, 346.9)
            (60, 329.8)
            (80, 312.9)
            (100, 167.1)
        }%
        }%
    }
    \vspace{-0.5em}
    \subfloat[wenige mittellange Testcases: $T$=10, $I$=1.000.000]{%
        \resizebox{0.45\textwidth}{!}{%
        \benchmarkplot{
            (0, 475.6)
            (20, 471.9)
            (40, 485.6)
            (60, 483.0)
            (80, 488.4)
            (100, 487.2)
        }{
            (0, 546.8)
            (20, 497.2)
            (40, 446.8)
            (60, 396.5)
            (80, 341.5)
            (100, 163.1)
        }%
        }%
    }
    \hfill
    \subfloat[viele mittellange Testcases: $T$=1.000, $I$=10.000]{%
        \resizebox{0.45\textwidth}{!}{%
        \benchmarkplot{
            (0, 841.5)
            (20, 831.4)
            (40, 815.7)
            (60, 825.0)
            (80, 830.2)
            (100, 838.8)
        }{
            (0, 1558.4)
            (20, 1307.2)
            (40, 1087.1)
            (60, 829.2)
            (80, 591.2)
            (100, 202.8)
        }%
        }%
    }
    \vspace{-0.5em}
    \subfloat[wenige langsame Testcases: $T$=10, $I$=10.000.000]{%
        \resizebox{0.45\textwidth}{!}{%
        \benchmarkplot{
            (0, 3237.9)
            (20, 3417.2)
            (40, 3410.0)
            (60, 3335.6)
            (80, 3217.6)
            (100, 3075.6)
        }{
            (0, 2783.8)
            (20, 2101.1)
            (40, 1871.0)
            (60, 1236.2)
            (80, 743.7)
            (100, 164.0)
        }%
        }%
    }
    \hfill
    \subfloat[viele langsame Testcases: $T$=10.000, $I$=10.000]{%
        \resizebox{0.45\textwidth}{!}{%
        \benchmarkplot{
            (0, 7036.4)
            (20, 7137.5)
            (40, 7154.3)
            (60, 7017.1)
            (80, 7087.7)
            (100, 7054.8)
        }{
            (0, 21107.8)
            (20, 17356.5)
            (40, 13287.1)
            (60, 9234.9)
            (80, 5136.3)
            (100, 830.7)
        }%
        }%
    }
    \caption{Ergebnisse der Benchmark. Sämtliche dieser Übersicht zugrunde liegenden Daten können in dem Repository der Benchmark (siehe \chapref{Links}) eingesehen und überprüft werden.}
    \label{fig:sechs-grafiken}
\end{figure}

\begin{figure}[htbp]
    \centering
    \subfloat[wenige kurze Testcases: $T$=10, $I$=100]{%
        \resizebox{0.45\textwidth}{!}{%
        \benchmarkplotOverhead{
            (0, 122.0)
            (20, 122.4)
            (40, 124.5)
            (60, 129.2)
            (80, 120.6)
            (100, 121.6)}
            {
                (0, 3.6)
                (20, 3.6)
                (40, 3.6)
                (60, 3.6)
                (80, 3.5)
                (100, 3.6)
            }{
                (0, 210.2)
                (20, 138.8)
                (40, 138.9)
                (60, 134.9)
                (80, 134.9)
                (100, 9.6)
            }{
                (0, 39.6)
                (20, 41.1)
                (40, 41.6)
                (60, 41.5)
                (80, 41.1)
                (100, 39.3)
            }%
        }%
    }
    \hfill
    \subfloat[viele kurze Testcases: $T$=100, $I$=10]{%
        \resizebox{0.45\textwidth}{!}{%
        \benchmarkplotOverhead{
            (0, 161.8)
            (20, 165.9)
            (40, 162.5)
            (60, 163.0)
            (80, 161.2)
            (100, 164.2)
        }{
            (0, 382.8)
            (20, 370.0)
            (40, 352.5)
            (60, 336.6)
            (80, 367.9)
            (100, 177.9)
        }{}{}%
        }%
    }
    \vspace{-0.5em}
    \subfloat[wenige mittellange Testcases: $T$=10, $I$=1.000.000]{%
        \resizebox{0.45\textwidth}{!}{%
        \benchmarkplotOverhead{
            (0, 456.6)
            (20, 427.6)
            (40, 372.0)
            (60, 443.3)
            (80, 357.4)
            (100, 363.5)
        }{
            (0, 553.0)
            (20, 523.3)
            (40, 472.8)
            (60, 400.1)
            (80, 347.3)
            (100, 175.0)
        }{}{}%
        }%
    }
    \hfill
    \subfloat[viele mittellange Testcases: $T$=1.000, $I$=10.000]{%
        \resizebox{0.45\textwidth}{!}{%
        \benchmarkplotOverhead{
            (0, 783.7)
            (20, 789.6)
            (40, 778.1)
            (60, 843.8)
            (80, 803.6)
            (100, 787.3)
        }{
            (0, 1599.1)
            (20, 1318.1)
            (40, 1063.7)
            (60, 830.1)
            (80, 596.4)
            (100, 215.2)
        }{}{}%
        }%
    }
    \vspace{-0.5em}
    \subfloat[wenige langsame Testcases: $T$=10, $I$=10.000.000]{%
        \resizebox{0.45\textwidth}{!}{%
        \benchmarkplotOverhead{
            (0, 2580.0)
            (20, 2418.9)
            (40, 2739.4)
            (60, 2611.9)
            (80, 2666.4)
            (100, 2378.4)
        }{
            (0, 3070.7)
            (20, 2237.5)
            (40, 1755.5)
            (60, 1239.1)
            (80, 783.0)
            (100, 228.3)
        }{}{}%
        }%
    }
    \hfill
    \subfloat[viele langsame Testcases: $T$=10.000, $I$=10.000]{%
        \resizebox{0.45\textwidth}{!}{%
        \benchmarkplotOverhead{
            (0, 7118.8)
            (20, 7074.7)
            (40, 7068.8)
            (60, 7080.6)
            (80, 7115.5)
            (100, 7116.6)
        }{
            (0, 21777.5)
            (20, 17293.3)
            (40, 13253.1)
            (60, 9229.9)
            (80, 5158.9)
            (100, 831.9)
        }{}{}%
        }%
    }
    \caption{Detailansicht zu den Laufzeiten einzelner Komponenten des Verfahrens. Sämtliche dieser Übersicht zugrunde liegenden Daten können in dem Repository der Benchmark (siehe \chapref{Links}) eingesehen und überprüft werden.}
    \label{fig:benchmark_overhead}
\end{figure}

\section{Zusammenfassung der Ergebnisse}
Aus \figref{fig:benchmark_evaluation} ist ersichtlich, dass sich mit der Anwedung der Verfahrens durchaus eine Verrignerung der Laufzeit eines Testdurchlaufs erreichen lässt. Durch die Verwendung der verschiedenen Komponenten des Verfahrens und der notwendigen Kommunikation dieser untereinander entsteht zwar ein gewisser Overhead im Vergleich zu einem konventionellen Testdurchlauf, dieser nimmt jedoch mit zunehmender Abdeckung von Testcases in \texttt{archive} nahezu konstant ab.\\
Jedoch decken die Graphen (2), (4) und (6) einen Makel auf: Bei einer hohen Anzahl an Testcases, die ausgeführt werden müssen, ist die Laufzeit des Verfahrens zum Teil deutlich höher als ein konventioneller Durchlauf. \figref{fig:benchmark_overhead} liefert die Erklärung zu diesem Verhalten: Während die Laufzeiten für die \texttt{testDiscovery}, den \texttt{testAuditor} und für sonstigen Overhead, der durch das ausführende Script entsteht, innerhalb einer Konfiguration der Benchmark nahezu konstant bleiben, ist die Laufzeit der \texttt{testExecution} stark von der vorhandenen Menge an Testcases abhängig. Vergleicht man einzelne Graphen von der gleichen Konfiguration aus \figref{fig:benchmark_evaluation} und \figref{fig:benchmark_overhead} miteinander, so ist ersichtlich, dass die \texttt{testExecution} bei auch gleicher Menge an Testcases deutlich langsamer ist als ein konventioneller Testlauf. Dies lässt sich durch das verwendete Testframework und die Funktionsweise der \texttt{testExecution} erklären:\\
Die \texttt{testExecution} bekommt von \texttt{testAuditor} eine Menge an Namen von Testcases, die ausgeführt werden sollen. In dieser Benchmark wird diese Menge von der \texttt{testExecution} an pytest übergeben, was bei einer hohen Anzahl von Testcases eine sehr teure Operation werden kann. Außerdem ist pytest nicht effizient auf diese Art von Benutzung ausgelegt, sodass die Laufzeit bei dieser Art von Aufruf mit vielen Testcases deutlich schlechter als ein kompletter Neudurchlauf ist. Der Grad der Auswirkung dieser Implementationsweise ist also auch abhängig vom verwendeten Testframework.\\
Die beste Performance zeigt das Verfahren jedoch in Fällen wie in Graph (5) demonstriert:\\
Bei einer geringen Anzahl fällt die langsame Implementation der \texttt{testExecution} weniger ins Gewicht, und bei einer langen Laufzeit der einzelnen Testcases kann durch das Vermeiden des Ausführens von Testcases, für die bereits ein Ergebnis vorliegt, eine signifikate Verringerung der Gesamtlaufzeit des Testdurchlaufs erzielt werden. So performt das Verfahren in (5) ab einer Abdeckung der Testcases von etwa 20\% bereits besser als der konventionelle Durchlauf.\\
Je mehr Testcases bereits ausgeführt wurden, desto bessere Laufzeiten erzielt das Verfahren. Selbst bei Konfigurationen mit vielen Testcases performt das Verfahren ab einer Abdeckung der Testcases durch \texttt{archive} von etwa 80\% besser als ein konventioneller Durchlauf. Eine Ausnahmen bilden Konfigurationen mit Testcases mit sehr kurzer Laufzeit: Hier ist der Laufzeitgewinn pro gespartem Testcase gering, und der Overhead überwiegt die Laufzeit eines konventionellen Durchlaufs, wie in (1) und (2) ersichtlich. 
